\documentclass[10pt,a4paper]{amsart}
\usepackage[margin=19mm]{geometry}
\usepackage{amsmath,amssymb,mathtools}
\usepackage[scr=rsfs,cal=euler]{mathalfa}
\usepackage{xfrac}
\usepackage[unicode,hidelinks,bookmarks=false]{hyperref}
\newcommand{\ie}{i.e.\ }
\newcommand{\cf}{cf.\ }
\newcommand{\resp}{respectively\ }
\newcommand{\Subsection}{\subsection}
\providecommand{\nobreakdash}{\nobreak\hbox{-}}
\setlength{\emergencystretch}{2em}
\multlinegap0pt
\usepackage{amssymb}
\usepackage{mathtools}
\usepackage{xcolor}
\newtheorem{res}{Result}
\newtheorem{lm}{Lemma}
\newtheorem{sideways}{Sideways}
\newtheorem{prop}{Proposition}
\usepackage{cleveref}
\crefname{res}{Result}{Results}
\crefname{lm}{Lemma}{Lemmas}
\crefname{sideways}{Sideways}{Sidewayss}
\crefname{prop}{Proposition}{Propositions}
 \DeclareMathOperator{\dds}{dds}
 \newcommand{\eps}{\varepsilon}
 \newcommand{\mq}{\mathcal Q}
 \newcommand{\ms}{\mathcal S}
 \newcommand{\one}{\mathbf 1}
 \newcommand{\pp}{\mathcal P}
 \newcommand{\sett}[2]{ \left\{ #1 \, \, || \, \, #2 \right \} }
 \newcommand{\vspan}[1]{\left \langle #1 \right \rangle}
 \newcommand{\vspann}[2]{\left \langle #1 \, \| \, #2 \right \rangle}
\title{Association schemes and orthogonality graphs on anisotropic points of polar spaces}
\author{Sam Adriaensen, Maarten De Boeck}
\date{Source: arXiv:2402.05055v2 (CC BY 4.0)}
\begin{document}
\maketitle

\noindent\emph{Source and licence.} Association schemes and orthogonality graphs on anisotropic points of polar spaces. Version arXiv:2402.05055v2 (CC BY 4.0), distributed by arXiv under the Creative Commons Attribution 4.0 International licence, http://creativecommons.org/licenses/by/4.0/, as recorded in the arXiv licence metadata for this exact version and shown on the abstract page https://arxiv.org/abs/2402.05055v2. Full licence notice: \texttt{licenses/arxiv-2402.05055-CC-BY-4.0.txt}.\par\smallskip

\noindent\textit{Editorial scope.} Counted unit: the article's prop Prop:RDisjointHoffman (source0.tex lines 1609--1614). The article's complete original proof of it is delivered with it (source0.tex lines 1616--1647, one \texttt{proof} environment of 32 lines). No mathematics is written, repaired or re-derived: the statement and the proof are the article's own lines, in the article's own notation, and no retained line is substituted or re-flowed.  Retained uncounted as the source-local context the proof needs: lines 582--591; lines 601--610; lines 1469--1541.  Nothing was omitted from any retained span, so the package carries no omission record.  No retained line contains a citation, so the wrapper carries no bibliography and no bibliography entry.  No retained line contains a figure environment, an image-inclusion command or drawing code, so no image is delivered and the package declares no asset dependency.  Editorial formatting only, in addition: an AMS \texttt{amsart} preamble that restates the article's own declarations and macros used by the retained lines, namely the environments \texttt{res}, \texttt{lm}, \texttt{sideways}, \texttt{prop} on separate counters, together with the standard packages those lines need.  The retained environments print in this document as Res 1, Lm 1, Sideways 1, Proposition 1 in order of appearance, whereas the article prints the same items with the numbering of its own full document.  The complete native source of the article as distributed in the arXiv e-print for this exact version is bundled unchanged as \texttt{sources/154191/source0.tex}.
\begin{res}[Weighted Hoffman's ratio bound]
 \label{Res:Hoffman}
 Let $(\pp,R)$ be a graph and $A$ a compatible symmetric matrix.
 Suppose that $A \one = n \one$ and that $\tau < n$ is the smallest eigenvalue of $A$.
 If $S$ is a coclique in $(\pp,R)$, then
 \[
  |S| \leq \frac{|\pp|}{\frac{n}{-\tau}+1}.
 \]
 In case of equality, $\chi_S$ is contained in the span of $\one$ and the $\tau$-eigenspace of $A$.
\end{res}

\begin{lm}
 \label{Lm:SpanningOrbit}
 Let $S$ be a non-empty subset of $\pp$.
 Let $G$ be a group of automorphisms such that for any $X,Y \in \pp$, the number $|\sett{g \in G}{X,Y \in S^g}|$ only depends on the relation containing $(X,Y)$.
 Then $\vspan{\chi_{S^g} \, \| \, g \in G} = \vspan{V_j \, \| \, j \in \dds(S) \cup \{0\}}$.
 % \begin{enumerate}
 %  \item \label{Lm:SpanningOrbit:Span} 
 %  \item \label{Lm:SpanningOrbit:Disjoint} The dual degree sets of $S$ and $T$ are disjoint if and only if $|S^g \cap T|$ is equal for all $g \in G$.
 % \end{enumerate}
\end{lm}

\begin{table}
 \centering
\begin{minipage}{.4\textwidth}
 \begin{center}
\begin{sideways}
 $\displaystyle \mathbf P = \begin{pmatrix}
 1 & q^{n-1}-1 
   & \frac14 q^{\frac{n-1}2}\left(q^{\frac{n-1}2} + \eps\right)(q-3) 
   & \frac14 q^{\frac{n-1}2}\left(q^{\frac{n-1}2} - \eps\right)(q-1) 
   & \frac14 q^{\frac{n-1}2}\left(q^{\frac{n-1}2} + \eps\right)(q-1) 
   & \frac14 q^{\frac{n-1}2}\left(q^{\frac{n-1}2} - \eps\right)(q+1) \\[1em]
 1 & q^{n-1}-1 
   & \frac14 q^{\frac{n-1}2}\left(q^{\frac{n-1}2} + \eps\right)(q-3) 
   & \frac14 q^{\frac{n-1}2}\left(q^{\frac{n-1}2} - \eps\right)(q-1) 
   & -\frac14 q^{\frac{n-1}2}\left(q^{\frac{n-1}2} + \eps\right)(q-1) 
   & -\frac14 q^{\frac{n-1}2}\left(q^{\frac{n-1}2} - \eps\right)(q+1) \\[1em]
 1 & \eps q^{\frac{n-3}2}-1 
   & \frac14 \eps q^{\frac{n-3}2}(q+1)(q-3)
   & -\frac14 \eps q^{\frac{n-3}2}(q-1)^2
   & \frac14 \eps q^{\frac{n-3}2}(q-1)(q+1)
   & - \frac14 \eps q^{\frac{n-3}2}(q-1)(q+1) \\[1em]
 1 & \eps q^{\frac{n-3}2}-1 
   & \frac14 \eps q^{\frac{n-3}2}(q+1)(q-3)
   & -\frac14 \eps q^{\frac{n-3}2}(q-1)^2
   & -\frac14 \eps q^{\frac{n-3}2}(q-1)(q+1)
   & \frac14 \eps q^{\frac{n-3}2}(q-1)(q+1) \\[1em]
 1 & \eps q^{\frac{n-1}2} - 1 & -\eps q^{\frac{n-1}2} & 0 & 0 & 0 \\[1em]
 1 & -\eps q^{\frac{n-1}2} - 1 & 0 & \eps q^{\frac{n-1}2} & 0 & 0
 \end{pmatrix}$
\end{sideways}
 \end{center}
\end{minipage}
\begin{minipage}{.4\textwidth}
 \begin{center}
\begin{sideways}
 $\displaystyle \mathbf Q = \begin{pmatrix}
 1 & 1 
  & q^2\frac{q^{n-1}-1}{q^2-1} 
  & q^2\frac{q^{n-1}-1}{q^2-1} 
  & \frac{q-3}{2(q-1)} \left( q^{\frac{n-1}2}+\eps\right) \left(q^{\frac{n+1}2}-\eps \right) 
  & \frac{q-1}{2(q+1)} \left( q^{\frac{n-1}2}-\eps\right) \left(q^{\frac{n+1}2}-\eps \right) \\[1em]
 1 & 1 
  & q^2 \frac{\eps q^{\frac{n-3}2}-1 }{ q^2-1 }
  & q^2 \frac{\eps q^{\frac{n-3}2}-1 }{ q^2-1 }
  & \eps \frac{q-3}{2(q-1)} \left( q^{\frac{n+1}2} - \eps \right)
  & -\eps \frac{q-1}{2(q+1)} \left( q^\frac{n+1}2 - \eps \right) \\[1em]
 1 & 1 
  & \eps \frac q{q-1} \left( q^{\frac{n-1}2} - \eps \right)
  & \eps \frac q{q-1} \left( q^{\frac{n-1}2} - \eps \right)
  & - \eps \frac{2}{q-1} \left( q^\frac{n+1}2 - \eps \right)
  & 0 \\[1em]
 1 & 1 
  & -\eps \frac q{q+1} \left( q^{\frac{n-1}2} + \eps \right)
  & -\eps \frac q{q+1} \left( q^{\frac{n-1}2} + \eps \right)
  & 0
  & \eps \frac2{q+1} \left( q^\frac{n+1}2 - \eps \right) \\[1em]
 1 & -1 
  & \eps \frac q{q-1} \left( q^{\frac{n-1}2} - \eps \right)
  & - \eps\frac q{q-1} \left( q^{\frac{n-1}2} - \eps \right) 
  & 0
  & 0 \\[1em]
 1 & -1 
  & - \eps \frac q{q+1} \left( q^{\frac{n-1}2} + \eps \right)
  & \eps \frac q{q+1} \left( q^{\frac{n-1}2} + \eps \right)
  & 0
  & 0
 \end{pmatrix}$
\end{sideways}
 \end{center}
\end{minipage}
 \caption{The matrices of eigenvalues and dual eigenvalues of the association scheme.}
 \label{Tab:Matrices}
\end{table}

\begin{prop}
 \label{Prop:RDisjointHoffman}
 Let $\ms_2$ denote the set of totally isotropic $\frac{n-3}2$-spaces.
 For every $\pi \in \ms_2$, $\pi^\perp \cap \pp$ is a weighted Hoffman coclique for the relation $R_\frac{5+\eps}2 \cup R_\frac{9+\eps}2$ of size $q^\frac{n-1}2(q-\eps)$.
 Moreover, $\vspann{\chi_{\pi^\perp}}{\pi \in \ms_2} = V_0 \oplus V_2$ and $\vspann{\chi_{\pi^\perp}^{S_q}}{\pi \in \ms_2} = V_1 \oplus V_3$.
\end{prop}

\begin{proof}
 If $\pi \in \ms_2$, then $\pi^\perp$ intersects $\mq^\eps(n,q)$ in a cone $\pi \mq^\eps(1,q)$.
 This implies that $\pi^\perp$ contains no lines intersecting $\mq^\eps(n,q)$ in exactly $1-\eps$ points, and hence that the anisotropic points of $\pi^\perp$ are a coclique for relations $R_{\frac{5+\eps}2}$ and $R_{\frac{9+\eps}{2}}$.
 The number of anisotropic points in $\pi^\perp$ equals
 \[
  \theta_{\frac{n+1}2}(q) - |\pi \mq^\eps(1,q)|
  = \theta_{\frac{n+1}2}(q) - \left( \theta_{\frac{n-3}2}(q) + (1+\eps) q^\frac{n-1}2 \right)
  = (q-\eps) q^\frac{n-1}2.
 \]
 We will check that this matches a weighted Hoffman's ratio bound.
 Note that any linear combination of $A_{\frac{5+\eps}2}$ and $A_{\frac{9+\eps}2}$ satisfies the conditions of \Cref{Res:Hoffman} with respect to the relevant graph.
 The eigenvalues of this linear combination can be easily obtained from the matrix $\mathbf P$ from \Cref{Tab:Matrices}.

 When making a linear combination of $A_{\frac{5+\eps}2}$ and $A_{\frac{9+\eps}2}$, we want it to have its smallest eigenvalue on $V_2$.
 In particular, we need to avoid that the smallest eigenvalue corresponds to $V_1$.
 Therefore, we take the linear combination $A = (q+\eps) A_{\frac{5+\eps}2} + (q-2+\eps) A_{\frac{9+\eps}2}$.
 Then $\one$ is an eigenvector of $A$ with eigenvalue
 \[
  (q+\eps) \mathbf P\left(0,\frac{5+\eps}2\right) + (q-2+\eps) \mathbf P\left(0,\frac{9+\eps}2\right)
  = \frac 12 q^\frac{n-1}2 \left( q^\frac{n-1}2 - 1 \right) (q+\eps) (q-2+\eps)
 \]
 and $A$ has its smallest eigenvalue
 \[
  (q+\eps) \mathbf P\left(2,\frac{5+\eps}2\right) + (q-2+\eps) \mathbf P\left(2,\frac{9+\eps}2\right)
  = - \frac 12 q^\frac{n-3}2 (q-\eps) (q+\eps) (q-2+\eps)
 \]
 on eigenspace $V_2$.
 Plugging this into \Cref{Res:Hoffman} yields a bound $(q-\eps) q^\frac{n-1}2$ on cocliques for $R_{\frac{5+\eps}2} \cup R_{\frac{9+\eps}2}$, which is achieved by the set of anisotropic points in a space $\pi^\perp$ with $\pi \in \ms$.
 Then $V_0 \oplus V_2 = \vspann{\chi_{\pi^\perp}}{\pi \in \ms_2}$ by \Cref{Lm:SpanningOrbit}.
 Since the bipartite graph $(\pp,R_4)$ has eigenvalues with opposite sign on $V_0$ and $V_1$, and on $V_2$ and $V_3$, we know that $V_1 = \sett{v^{S_q}}{v \in V_0}$ and $V_3 = \sett{v^{S_q}}{v \in V_2}$.
 Thus, $V_1 \oplus V_3 = \vspann{\chi_{\pi^\perp}^{S_q}}{\pi \in \ms_2}$.
\end{proof}

\end{document}
